\documentclass[11pt]{article}

\usepackage[margin=1in]{geometry}
\usepackage{amsmath,amssymb,amsthm}
\usepackage{microtype}

\newtheorem{proposition}{Proposition}
\newtheorem{example}{Counterexample}
\newtheorem{remark}{Remark}

\title{Decision-Value Acquisition for Runtime Safety Monitoring}
\author{}
\date{}

\begin{document}
\maketitle

\section{Setup}

Let \(Y\) denote the unknown safety label, \(H\) all information available
before an optional monitor is acquired, and \(Z\) the output of that monitor.
A downstream action \(a\in\mathcal{A}\) incurs loss \(\ell(a,Y)\).

For a realized cheap-information state \(H=h\), define the conditional Bayes
risk without the optional monitor by
\[
r(h)=\inf_{a\in\mathcal{A}}
\mathbb{E}\!\left[\ell(a,Y)\mid H=h\right].
\]

After acquiring the optional monitor and observing \(Z=z\), define
\[
r(h,z)=\inf_{a\in\mathcal{A}}
\mathbb{E}\!\left[\ell(a,Y)\mid H=h,Z=z\right].
\]

Let \(c(h)\geq 0\) be the acquisition cost. The conditional decision value is
\[
\Delta(h)=r(h)-\mathbb{E}\!\left[r(h,Z)\mid H=h\right].
\]

Thus, \(\Delta(h)\) is the expected reduction in optimal downstream decision
loss before acquisition cost is charged.

\section{Optimal two-monitor acquisition rule}

\begin{proposition}[Decision-value acquisition rule]
For a fixed cheap-information state \(h\), acquiring the optional monitor is
strictly optimal exactly when
\[
\Delta(h)>c(h).
\]
If \(\Delta(h)=c(h)\), acquisition and non-acquisition have equal conditional
risk.
\end{proposition}

\begin{proof}
Without acquisition, the minimum conditional risk is
\[
R_{\mathrm{skip}}(h)=r(h).
\]

With acquisition, the cost is paid and the action is optimized after observing
\(Z\), giving
\[
R_{\mathrm{acquire}}(h)
=
c(h)+\mathbb{E}\!\left[r(h,Z)\mid H=h\right].
\]

Acquisition is strictly preferable if and only if
\[
R_{\mathrm{acquire}}(h)<R_{\mathrm{skip}}(h).
\]
Substitution yields
\[
c(h)+\mathbb{E}\!\left[r(h,Z)\mid H=h\right]<r(h).
\]
Rearranging gives
\[
r(h)-\mathbb{E}\!\left[r(h,Z)\mid H=h\right]>c(h),
\]
which is exactly
\[
\Delta(h)>c(h).
\]
Equality gives equal conditional risk; the reverse inequality makes skipping
optimal.
\end{proof}

\begin{remark}[Value of information before cost]
A decision rule that observes \(Z\) may always ignore it. Therefore,
\[
\mathbb{E}\!\left[r(h,Z)\mid H=h\right]\leq r(h),
\]
and hence \(\Delta(h)\geq 0\). Acquisition is worthwhile only when this value
exceeds its cost.
\end{remark}

\section{Counterexample to uncertainty routing}

Consider binary \(Y\in\{0,1\}\), binary actions, and zero-one loss
\(\ell(a,Y)=\mathbf{1}\{a\neq Y\}\). There are two equally likely contexts,
\(C=A\) and \(C=B\). The cheap information includes context, but its class
posterior is identical:
\[
\Pr(Y=1\mid H=h_A)
=
\Pr(Y=1\mid H=h_B)
=
\frac{1}{2}.
\]
Therefore,
\[
r(h_A)=r(h_B)=\frac{1}{2},
\]
so the two examples have identical cheap-model uncertainty.

\begin{example}[Identical uncertainty, different monitor quality]
In context \(A\), the optional monitor is a binary symmetric signal with
accuracy \(0.9\):
\[
\Pr(Z=Y\mid H=h_A)=0.9.
\]
The Bayes action after observing \(Z\) is \(a=Z\), whose risk is \(0.1\).
Therefore,
\[
\Delta(h_A)=0.5-0.1=0.4.
\]

In context \(B\), the optional monitor is independent of \(Y\):
\[
\Pr(Z=1\mid Y,H=h_B)=0.5.
\]
The posterior remains \(0.5\) after observing \(Z\), so
\[
\Delta(h_B)=0.5-0.5=0.
\]

For every cost \(c\) satisfying \(0<c<0.4\), the optimal rule acquires in
context \(A\) and skips in context \(B\), although cheap uncertainty is
identical.
\end{example}

At a matched acquisition rate of \(50\%\), a decision-value policy acquires
every context-\(A\) example and no context-\(B\) example. Its average downstream
decision loss is
\[
\frac{1}{2}(0.1)+\frac{1}{2}(0.5)=0.3.
\]

An uncertainty policy cannot distinguish the two contexts. With
context-independent random tie-breaking, it acquires half of each context.
Its average downstream loss is
\[
\frac{1}{2}\left[\frac{1}{2}(0.1)+\frac{1}{2}(0.5)\right]
+
\frac{1}{2}(0.5)
=
0.4.
\]

The policies have the same acquisition rate and therefore the same expected
acquisition cost, but uncertainty routing has \(0.1\) greater decision loss.
It is therefore suboptimal.

\section{Implication}

The routing target should be monitor-specific expected decision improvement,
not harmfulness probability or uncertainty alone. The next empirical task is
to estimate \(\Delta(h)\) and compare decision-value, uncertainty, random, and
oracle acquisition at identical budgets.

\end{document}
